% Level 05 - Split policy
% Question answered: which data trains, which tunes, and which is allowed to produce a benchmark number?
% Facts: project notes (synthetic: 1,000,000 train recipes, 20,000 validation, test skipped; fine-tune real train/val/test;
% final test is real test only, also report the zero-shot pretrained model); report (synthetic never in the benchmark; splits grouped by
% event, site and date; hashed manifest; optional controlled-release test set).
\documentclass[tikz,border=10pt]{standalone}
% =============================================================================
% MethaneWatch diagram style  (shared by every diagram under diagrams/)
% -----------------------------------------------------------------------------
% Why one shared file: every diagram must look like part of the same set, and a
% colour change (for example a layer re-coloured) should be one edit, not 33.
%
% Usage, at the top of any diagram:
%   \documentclass[tikz,border=10pt]{standalone}
%   % =============================================================================
% MethaneWatch diagram style  (shared by every diagram under diagrams/)
% -----------------------------------------------------------------------------
% Why one shared file: every diagram must look like part of the same set, and a
% colour change (for example a layer re-coloured) should be one edit, not 33.
%
% Usage, at the top of any diagram:
%   \documentclass[tikz,border=10pt]{standalone}
%   % =============================================================================
% MethaneWatch diagram style  (shared by every diagram under diagrams/)
% -----------------------------------------------------------------------------
% Why one shared file: every diagram must look like part of the same set, and a
% colour change (for example a layer re-coloured) should be one edit, not 33.
%
% Usage, at the top of any diagram:
%   \documentclass[tikz,border=10pt]{standalone}
%   \input{../_common/mwstyle.tex}
%   \begin{document} ... \end{document}
%
% Colour meaning is fixed across the whole set (see diagrams/README.md). Two colours plus grey:
%   blue = MethaneWatch's own parts (application, data, models, science, services)
%   red  = evidence, security, audit and hard limits (the things that must never be bypassed)
%   grey = people, cloud and infrastructure, external systems, notes
% The older names green, orange, purple, teal and gold are kept as aliases so that no diagram
% source needs to change; they now resolve to blue (green, orange, purple) or grey (teal, gold).
% =============================================================================
\usepackage[T1]{fontenc}
\usepackage[scaled=0.95]{helvet}
\renewcommand{\familydefault}{\sfdefault}
\hyphenpenalty=10000 \exhyphenpenalty=10000 % never hyphenate words inside diagrams
\usepackage{pifont}
\usepackage{amssymb}
\usepackage{array}
\usepackage{booktabs}
\usetikzlibrary{positioning,arrows.meta,shapes.geometric,shapes.multipart,%
  shapes.misc,fit,backgrounds,calc,decorations.pathreplacing,matrix,patterns}

% ---- palette ----------------------------------------------------------------
\definecolor{mwink}{HTML}{1F2933}
\definecolor{mwgrey}{HTML}{6B7785}
\definecolor{mwline}{HTML}{4A5560}
\definecolor{mwlight}{HTML}{F2F4F7}
\definecolor{mwblue}{HTML}{1F6FB2}   \definecolor{mwblueL}{HTML}{E3EFFA}
\definecolor{mwred}{HTML}{B5483F}    \definecolor{mwredL}{HTML}{F8E3E0}
% aliases (see note above)
\colorlet{mwgreen}{mwblue}   \colorlet{mwgreenL}{mwblueL}
\colorlet{mworange}{mwblue}  \colorlet{mworangeL}{mwblueL}
\colorlet{mwpurple}{mwblue}  \colorlet{mwpurpleL}{mwblueL}
\colorlet{mwteal}{mwgrey}    \colorlet{mwtealL}{mwlight}
\colorlet{mwgold}{mwgrey}    \colorlet{mwgoldL}{mwlight}

% ---- node and edge styles ---------------------------------------------------
\tikzset{
  >={Stealth[length=2.2mm]},
  every node/.style={font=\sffamily\small,text=mwink},
  box/.style={rectangle,rounded corners=2.5pt,draw=mwline,line width=0.7pt,
    fill=white,align=center,inner sep=5pt,minimum height=7mm},
  blue/.style={draw=mwblue,fill=mwblueL},
  green/.style={draw=mwgreen,fill=mwgreenL},
  orange/.style={draw=mworange,fill=mworangeL},
  purple/.style={draw=mwpurple,fill=mwpurpleL},
  red/.style={draw=mwred,fill=mwredL},
  teal/.style={draw=mwteal,fill=mwtealL},
  gold/.style={draw=mwgold,fill=mwgoldL},
  grey/.style={draw=mwgrey,fill=mwlight},
  strong/.style={line width=1.6pt},
  ext/.style={box,dashed,draw=mwgrey,fill=mwlight},
  actor/.style={box,gold,rounded corners=8pt,minimum height=9mm},
  store/.style={cylinder,shape border rotate=90,aspect=0.22,draw=mwgreen,fill=mwgreenL,
    line width=0.7pt,align=center,inner sep=3pt,minimum height=11mm,minimum width=20mm},
  group/.style={rectangle,rounded corners=4pt,draw=mwline,dashed,line width=0.7pt,
    inner sep=7pt},
  glabel/.style={anchor=north west,font=\sffamily\bfseries\footnotesize,text=mwgrey,inner sep=2pt},
  note/.style={font=\sffamily\footnotesize,text=mwgrey,align=left,inner sep=1pt},
  tiny/.style={font=\sffamily\scriptsize},
  head/.style={font=\sffamily\bfseries\small},
  flow/.style={->,line width=0.8pt,draw=mwline},
  dflow/.style={->,line width=0.8pt,draw=mwline,dashed},
  bflow/.style={<->,line width=0.8pt,draw=mwline},
  lab/.style={font=\sffamily\scriptsize,text=mwgrey,fill=white,inner sep=1pt},
}

% ---- title / footer helpers -------------------------------------------------
% \dtitle{Finding or name}{One-line subtitle: what the diagram answers}
\newcommand{\dtitle}[2]{%
  \node[anchor=south west,font=\sffamily\bfseries\Large,text=mwink,inner sep=0pt] (dt)
    at ([yshift=12mm]current bounding box.north west) {#1};
  \node[anchor=north west,font=\sffamily\footnotesize,text=mwgrey,inner sep=0pt]
    at ([yshift=-1.2mm]dt.south west) {#2};}
% \dfoot{Source or reading note}
\newcommand{\dfoot}[2][14cm]{%
  \node[anchor=north west,note,text width=#1] at
    ([yshift=-4mm]current bounding box.south west) {#2};}
% in-node text helpers: \nm{Name} bold name, \tech{technology} small grey italic
\newcommand{\nm}[1]{{\bfseries\small #1}}
\newcommand{\tech}[1]{{\scriptsize\itshape\color{mwgrey} #1}}
% check mark used in matrices
\newcommand{\yes}{\ding{51}}

%   \begin{document} ... \end{document}
%
% Colour meaning is fixed across the whole set (see diagrams/README.md). Two colours plus grey:
%   blue = MethaneWatch's own parts (application, data, models, science, services)
%   red  = evidence, security, audit and hard limits (the things that must never be bypassed)
%   grey = people, cloud and infrastructure, external systems, notes
% The older names green, orange, purple, teal and gold are kept as aliases so that no diagram
% source needs to change; they now resolve to blue (green, orange, purple) or grey (teal, gold).
% =============================================================================
\usepackage[T1]{fontenc}
\usepackage[scaled=0.95]{helvet}
\renewcommand{\familydefault}{\sfdefault}
\hyphenpenalty=10000 \exhyphenpenalty=10000 % never hyphenate words inside diagrams
\usepackage{pifont}
\usepackage{amssymb}
\usepackage{array}
\usepackage{booktabs}
\usetikzlibrary{positioning,arrows.meta,shapes.geometric,shapes.multipart,%
  shapes.misc,fit,backgrounds,calc,decorations.pathreplacing,matrix,patterns}

% ---- palette ----------------------------------------------------------------
\definecolor{mwink}{HTML}{1F2933}
\definecolor{mwgrey}{HTML}{6B7785}
\definecolor{mwline}{HTML}{4A5560}
\definecolor{mwlight}{HTML}{F2F4F7}
\definecolor{mwblue}{HTML}{1F6FB2}   \definecolor{mwblueL}{HTML}{E3EFFA}
\definecolor{mwred}{HTML}{B5483F}    \definecolor{mwredL}{HTML}{F8E3E0}
% aliases (see note above)
\colorlet{mwgreen}{mwblue}   \colorlet{mwgreenL}{mwblueL}
\colorlet{mworange}{mwblue}  \colorlet{mworangeL}{mwblueL}
\colorlet{mwpurple}{mwblue}  \colorlet{mwpurpleL}{mwblueL}
\colorlet{mwteal}{mwgrey}    \colorlet{mwtealL}{mwlight}
\colorlet{mwgold}{mwgrey}    \colorlet{mwgoldL}{mwlight}

% ---- node and edge styles ---------------------------------------------------
\tikzset{
  >={Stealth[length=2.2mm]},
  every node/.style={font=\sffamily\small,text=mwink},
  box/.style={rectangle,rounded corners=2.5pt,draw=mwline,line width=0.7pt,
    fill=white,align=center,inner sep=5pt,minimum height=7mm},
  blue/.style={draw=mwblue,fill=mwblueL},
  green/.style={draw=mwgreen,fill=mwgreenL},
  orange/.style={draw=mworange,fill=mworangeL},
  purple/.style={draw=mwpurple,fill=mwpurpleL},
  red/.style={draw=mwred,fill=mwredL},
  teal/.style={draw=mwteal,fill=mwtealL},
  gold/.style={draw=mwgold,fill=mwgoldL},
  grey/.style={draw=mwgrey,fill=mwlight},
  strong/.style={line width=1.6pt},
  ext/.style={box,dashed,draw=mwgrey,fill=mwlight},
  actor/.style={box,gold,rounded corners=8pt,minimum height=9mm},
  store/.style={cylinder,shape border rotate=90,aspect=0.22,draw=mwgreen,fill=mwgreenL,
    line width=0.7pt,align=center,inner sep=3pt,minimum height=11mm,minimum width=20mm},
  group/.style={rectangle,rounded corners=4pt,draw=mwline,dashed,line width=0.7pt,
    inner sep=7pt},
  glabel/.style={anchor=north west,font=\sffamily\bfseries\footnotesize,text=mwgrey,inner sep=2pt},
  note/.style={font=\sffamily\footnotesize,text=mwgrey,align=left,inner sep=1pt},
  tiny/.style={font=\sffamily\scriptsize},
  head/.style={font=\sffamily\bfseries\small},
  flow/.style={->,line width=0.8pt,draw=mwline},
  dflow/.style={->,line width=0.8pt,draw=mwline,dashed},
  bflow/.style={<->,line width=0.8pt,draw=mwline},
  lab/.style={font=\sffamily\scriptsize,text=mwgrey,fill=white,inner sep=1pt},
}

% ---- title / footer helpers -------------------------------------------------
% \dtitle{Finding or name}{One-line subtitle: what the diagram answers}
\newcommand{\dtitle}[2]{%
  \node[anchor=south west,font=\sffamily\bfseries\Large,text=mwink,inner sep=0pt] (dt)
    at ([yshift=12mm]current bounding box.north west) {#1};
  \node[anchor=north west,font=\sffamily\footnotesize,text=mwgrey,inner sep=0pt]
    at ([yshift=-1.2mm]dt.south west) {#2};}
% \dfoot{Source or reading note}
\newcommand{\dfoot}[2][14cm]{%
  \node[anchor=north west,note,text width=#1] at
    ([yshift=-4mm]current bounding box.south west) {#2};}
% in-node text helpers: \nm{Name} bold name, \tech{technology} small grey italic
\newcommand{\nm}[1]{{\bfseries\small #1}}
\newcommand{\tech}[1]{{\scriptsize\itshape\color{mwgrey} #1}}
% check mark used in matrices
\newcommand{\yes}{\ding{51}}

%   \begin{document} ... \end{document}
%
% Colour meaning is fixed across the whole set (see diagrams/README.md). Two colours plus grey:
%   blue = MethaneWatch's own parts (application, data, models, science, services)
%   red  = evidence, security, audit and hard limits (the things that must never be bypassed)
%   grey = people, cloud and infrastructure, external systems, notes
% The older names green, orange, purple, teal and gold are kept as aliases so that no diagram
% source needs to change; they now resolve to blue (green, orange, purple) or grey (teal, gold).
% =============================================================================
\usepackage[T1]{fontenc}
\usepackage[scaled=0.95]{helvet}
\renewcommand{\familydefault}{\sfdefault}
\hyphenpenalty=10000 \exhyphenpenalty=10000 % never hyphenate words inside diagrams
\usepackage{pifont}
\usepackage{amssymb}
\usepackage{array}
\usepackage{booktabs}
\usetikzlibrary{positioning,arrows.meta,shapes.geometric,shapes.multipart,%
  shapes.misc,fit,backgrounds,calc,decorations.pathreplacing,matrix,patterns}

% ---- palette ----------------------------------------------------------------
\definecolor{mwink}{HTML}{1F2933}
\definecolor{mwgrey}{HTML}{6B7785}
\definecolor{mwline}{HTML}{4A5560}
\definecolor{mwlight}{HTML}{F2F4F7}
\definecolor{mwblue}{HTML}{1F6FB2}   \definecolor{mwblueL}{HTML}{E3EFFA}
\definecolor{mwred}{HTML}{B5483F}    \definecolor{mwredL}{HTML}{F8E3E0}
% aliases (see note above)
\colorlet{mwgreen}{mwblue}   \colorlet{mwgreenL}{mwblueL}
\colorlet{mworange}{mwblue}  \colorlet{mworangeL}{mwblueL}
\colorlet{mwpurple}{mwblue}  \colorlet{mwpurpleL}{mwblueL}
\colorlet{mwteal}{mwgrey}    \colorlet{mwtealL}{mwlight}
\colorlet{mwgold}{mwgrey}    \colorlet{mwgoldL}{mwlight}

% ---- node and edge styles ---------------------------------------------------
\tikzset{
  >={Stealth[length=2.2mm]},
  every node/.style={font=\sffamily\small,text=mwink},
  box/.style={rectangle,rounded corners=2.5pt,draw=mwline,line width=0.7pt,
    fill=white,align=center,inner sep=5pt,minimum height=7mm},
  blue/.style={draw=mwblue,fill=mwblueL},
  green/.style={draw=mwgreen,fill=mwgreenL},
  orange/.style={draw=mworange,fill=mworangeL},
  purple/.style={draw=mwpurple,fill=mwpurpleL},
  red/.style={draw=mwred,fill=mwredL},
  teal/.style={draw=mwteal,fill=mwtealL},
  gold/.style={draw=mwgold,fill=mwgoldL},
  grey/.style={draw=mwgrey,fill=mwlight},
  strong/.style={line width=1.6pt},
  ext/.style={box,dashed,draw=mwgrey,fill=mwlight},
  actor/.style={box,gold,rounded corners=8pt,minimum height=9mm},
  store/.style={cylinder,shape border rotate=90,aspect=0.22,draw=mwgreen,fill=mwgreenL,
    line width=0.7pt,align=center,inner sep=3pt,minimum height=11mm,minimum width=20mm},
  group/.style={rectangle,rounded corners=4pt,draw=mwline,dashed,line width=0.7pt,
    inner sep=7pt},
  glabel/.style={anchor=north west,font=\sffamily\bfseries\footnotesize,text=mwgrey,inner sep=2pt},
  note/.style={font=\sffamily\footnotesize,text=mwgrey,align=left,inner sep=1pt},
  tiny/.style={font=\sffamily\scriptsize},
  head/.style={font=\sffamily\bfseries\small},
  flow/.style={->,line width=0.8pt,draw=mwline},
  dflow/.style={->,line width=0.8pt,draw=mwline,dashed},
  bflow/.style={<->,line width=0.8pt,draw=mwline},
  lab/.style={font=\sffamily\scriptsize,text=mwgrey,fill=white,inner sep=1pt},
}

% ---- title / footer helpers -------------------------------------------------
% \dtitle{Finding or name}{One-line subtitle: what the diagram answers}
\newcommand{\dtitle}[2]{%
  \node[anchor=south west,font=\sffamily\bfseries\Large,text=mwink,inner sep=0pt] (dt)
    at ([yshift=12mm]current bounding box.north west) {#1};
  \node[anchor=north west,font=\sffamily\footnotesize,text=mwgrey,inner sep=0pt]
    at ([yshift=-1.2mm]dt.south west) {#2};}
% \dfoot{Source or reading note}
\newcommand{\dfoot}[2][14cm]{%
  \node[anchor=north west,note,text width=#1] at
    ([yshift=-4mm]current bounding box.south west) {#2};}
% in-node text helpers: \nm{Name} bold name, \tech{technology} small grey italic
\newcommand{\nm}[1]{{\bfseries\small #1}}
\newcommand{\tech}[1]{{\scriptsize\itshape\color{mwgrey} #1}}
% check mark used in matrices
\newcommand{\yes}{\ding{51}}

\begin{document}
\begin{tikzpicture}[
  d/.style={box,text width=3.5cm,align=left,inner sep=4pt,minimum height=1.7cm,font=\sffamily\footnotesize}]

  \node[head,anchor=west] at (-0.2,2.4) {Source};
  \node[head,anchor=west] at (5.2,2.4) {Split};
  \node[head,anchor=west] at (10.6,2.4) {Used for};

  % row 1: synthetic
  \node[d,green] (s1) at (1.6,0) {\nm{Synthetic pretraining}\\ MethaneSET backgrounds + simulated plumes};
  \node[d,grey] (a1) at (6.9,1.1) {\nm{train}\\ 1,000,000 recipes};
  \node[d,grey] (a2) at (6.9,-1.1) {\nm{validation}\\ 20,000 recipes; test skipped};
  \node[d,orange,text width=4.2cm] (u1) at (12.5,0) {\nm{Pretraining and early stopping only}\\ Never reported as a result};

  % row 2: real
  \node[d,green] (s2) at (1.6,-4.4) {\nm{Real fine-tuning data}\\ MethaneSET fine-tune set, real scenes};
  \node[d,grey] (b1) at (6.9,-3.3) {\nm{train}\\ grouped by event, site, date};
  \node[d,grey] (b2) at (6.9,-5.5) {\nm{validation}\\ model selection};
  \node[d,red,strong] (b3) at (6.9,-7.7) {\nm{real test}\\ final scoring only};
  \node[d,orange,text width=4.2cm] (u2) at (12.5,-4.4) {\nm{Fine-tuning and model selection}};
  \node[d,red,strong,text width=4.2cm] (u3) at (12.5,-7.7) {\nm{The benchmark}\\ All models, incl. zero-shot pretrained, scored here};

  % optional
  \node[d,gold,dashed] (s3) at (1.6,-7.7) {\nm{Controlled-release data}\\\tech{optional, if access is granted}};

  \draw[flow] (s1) -- (a1.west); \draw[flow] (s1) -- (a2.west);
  \draw[flow] (a1.east) -- (u1.west |- a1.east) -- (u1.north west);
  \draw[flow] (a2.east) -- (u1.west |- a2.east) -- (u1.south west);
  \draw[flow] (s2) -- (b1.west); \draw[flow] (s2) -- (b2.west); \draw[flow] (s2.south east) -- (b3.north west);
  \draw[flow] (b1.east) -- (u2.north west); \draw[flow] (b2.east) -- (u2.south west);
  \draw[flow] (b3) -- (u3);
  \draw[dflow] (s3) -- (b3);

  % forbidden arrow
  \draw[red!70!black,line width=1.2pt,-{Stealth[length=3mm]}] (u1.east) -- ++(1.0,0) |- (u3.east);
  \node[cross out,draw=mwred,line width=1.6pt,minimum size=9pt] at (15.6,-3.8) {};
  \node[font=\sffamily\bfseries\footnotesize,text=mwred,anchor=west] at (15.8,-3.8) {never};

  \dtitle{Synthetic data teaches; only real data judges}
         {Level 05 - split policy. A benchmark number can come from the real test split only.}
  \dfoot[17cm]{Guards the code must enforce: a manifest listing every sample id with its split and a hash; grouping by event, site and date so neighbours cannot leak across splits; a test that fails if any synthetic id appears in the benchmark set.}
\end{tikzpicture}
\end{document}
